\documentclass{article}
\usepackage{amsmath,amssymb,amsthm}
\usepackage{algorithm}
\usepackage{algpseudocode}
\usepackage{geometry}
\geometry{margin=1in}
\newtheorem{theorem}{Theorem}
\newtheorem{lemma}{Lemma}
\newcommand{\E}{\mathbb{E}}
\newcommand{\R}{\mathbb{R}}
\newcommand{\norm}[1]{\left\lVert #1\right\rVert}
\begin{document}

\section*{Clipped Stochastic Gradient Descent (Clipped SGD)}

\section*{Mathematical Formulation}
Let $f:\R^{d}\rightarrow\R$ be the (possibly non‑convex) objective.  
At iteration $t$ we draw a stochastic gradient $g_t^{\text{raw}}$ satisfying the usual unbiasedness and variance assumptions (see Section~\ref{sec:assumptions}).  
The \emph{gradient clipping operator} with threshold $C>0$ is defined as  

\[
\operatorname{clip}(u;C)\;=\;
\frac{u}{\max\!\bigl(1,\;\norm{u}/C\bigr)}\;,
\qquad u\in\R^{d}.
\tag{1}
\]

Thus the clipped stochastic gradient is  

\[
g_t \;=\;\operatorname{clip}\bigl(g_t^{\text{raw}};C\bigr)
      \;=\;\frac{g_t^{\text{raw}}}{\max\!\bigl(1,\;\norm{g_t^{\text{raw}}}/C\bigr)} .
\tag{2}
\]

The update rule of Clipped SGD is  

\[
x_{t+1}=x_t-\eta_t g_t ,\qquad t=0,1,\dots ,T-1,
\tag{3}
\]

where $\{\eta_t\}_{t\ge 0}$ is a stepsize (learning‑rate) schedule.

\section*{Pseudocode}
\begin{algorithm}[H]
\caption{Clipped SGD}
\label{alg:clippedsgd}
\begin{algorithmic}[1]
\Require Initial point $x_0\in\R^{d}$, stepsizes $\{\eta_t\}$, clipping threshold $C>0$
\For{$t=0$ {\bf to} $T-1$}
    \State Sample a stochastic gradient $g_t^{\text{raw}}$ such that 
          $\E[g_t^{\text{raw}}\mid x_t]=\nabla f(x_t)$
    \State $g_t\gets g_t^{\text{raw}}/
          \max\bigl(1,\norm{g_t^{\text{raw}}}/C\bigr)$   \Comment{clipping}
    \State $x_{t+1}\gets x_t-\eta_t g_t$
\EndFor
\State \Return $x_T$ (or optionally the average $\bar x_T=\frac1T\sum_{t=0}^{T-1}x_t$)
\end{algorithmic}
\end{algorithm}

\section*{Dry Run Example}
Consider the one‑dimensional quadratic $f(w)=(w-3)^2$.  
Its gradient is $\nabla f(w)=2(w-3)$.  
Set $w_0=0$, stepsize $\eta=0.1$, clipping threshold $C=1$ and run three iterations.

\begin{center}
\begin{tabular}{c|c|c|c|c}
$t$ & $w_t$ & $\nabla f(w_t)$ & $g_t$ (clipped) & $f(w_t)$ \\ \hline
0 & $0$ & $-6$ & $g_0=\frac{-6}{\max(1,6/1)}=-1$ & $9$ \\
1 & $w_1=0-0.1(-1)=0.1$ & $-5.8$ & $g_1=\frac{-5.8}{\max(1,5.8)}=-1$ & $8.41$ \\
2 & $w_2=0.1-0.1(-1)=0.2$ & $-5.6$ & $g_2=-1$ & $7.84$ \\
3 & $w_3=0.2-0.1(-1)=0.3$ & $-5.4$ & $g_3=-1$ & $7.29$
\end{tabular}
\end{center}
The clipping forces every stochastic (here exact) gradient to have magnitude at most $C=1$, while the objective value steadily decreases.

\section*{Assumptions}
\label{sec:assumptions}
\begin{enumerate}
\item \textbf{Smoothness.} $f$ is $L$‑smooth: for all $x,y\in\R^{d}$,
\[
f(y)\le f(x)+\langle\nabla f(x),y-x\rangle+\frac{L}{2}\norm{y-x}^{2}.
\tag{A1}
\]
\item \textbf{Unbiased stochastic gradients.} For every $t$,
\[
\E\!\left[g_t^{\text{raw}}\mid x_t\right]=\nabla f(x_t).
\tag{A2}
\]
\item \textbf{Bounded variance.} There exists $\sigma^{2}\ge0$ such that
\[
\E\!\left[\norm{g_t^{\text{raw}}-\nabla f(x_t)}^{2}\mid x_t\right]\le\sigma^{2}.
\tag{A3}
\]
\item \textbf{Clipping threshold.} $C>0$ is a fixed constant. Consequently,
\[
\norm{g_t}\le C,\qquad\forall t.
\tag{A4}
\]
\item \textbf{Stepsizes.} The stepsizes are constant $\eta_t\equiv\eta$ satisfying
\[
0<\eta\le\frac{1}{L}.
\tag{A5}
\]
\end{enumerate}

\section*{Main Convergence Theorem}
\begin{theorem}[Convergence of Clipped SGD]\label{thm:main}
Under Assumptions (A1)–(A5), for any $T\ge1$ the iterates generated by Algorithm~\ref{alg:clippedsgd} satisfy  

\[
\frac{1}{T}\sum_{t=0}^{T-1}\E\!\bigl[\norm{\nabla f(x_t)}^{2}\bigr]
\;\le\;
\underbrace{\frac{2\bigl(f(x_0)-\inf f\bigr)}{\eta T}}_{(\text{optimization term})}
\;+\;
\underbrace{L\eta\sigma^{2}}_{(\text{variance term})}
\;+\;
\underbrace{\frac{2L\eta C^{2}}{T}}_{(\text{clipping bias term})}.
\tag{4}
\]
Consequently, choosing $\eta=\Theta\!\bigl(1/\sqrt{T}\bigr)$ yields  

\[
\frac{1}{T}\sum_{t=0}^{T-1}\E\!\bigl[\norm{\nabla f(x_t)}^{2}\bigr]=
O\!\Bigl(\frac{1}{\sqrt{T}}\Bigr).
\]
\end{theorem}

\section*{Theoretical Properties}
\begin{itemize}
\item \textbf{Stability.} Because $\norm{g_t}\le C$, the update (3) cannot produce arbitrarily large jumps; the iterates remain in a bounded region whenever $f$ has bounded level sets.
\item \textbf{Hyper‑parameter sensitivity.} The bound (4) shows three distinct contributions:
    \begin{enumerate}
        \item A $1/(\eta T)$ term – decreasing $\eta$ improves the optimization term but harms the variance term.
        \item A linear $\eta$ variance term – the usual trade‑off as in vanilla SGD.
        \item An extra $C^{2}\eta/T$ term caused by clipping; it vanishes as $T\to\infty$ for any fixed $C$.
    \end{enumerate}
    Hence, moderate $C$ (large enough not to clip most gradients) and stepsize $\eta\asymp 1/\sqrt{T}$ give the optimal $O(T^{-1/2})$ rate.
\item \textbf{Comparison to baselines.}  
    - \emph{Plain SGD} (no clipping) enjoys the same $O(T^{-1/2})$ rate under identical assumptions, but without the clipping bias term.  
    - \emph{Clipped SGD} adds robustness: when occasional gradients explode, the clipping term caps their norm, preventing divergence; the price is the extra $C^{2}\eta/T$ term, which is negligible asymptotically.
\end{itemize}

\section*{Discussion}
Clipping introduces a deterministic bias because $g_t\neq\nabla f(x_t)$ in general. The proof below carefully separates the bias from the stochastic noise and shows that the bias contributes only an $O(\eta C^{2}/T)$ additive term. This mirrors practical observations: clipping stabilises training while preserving the optimal $1/\sqrt{T}$ convergence order for smooth non‑convex objectives.

\section*{Preliminaries (Definitions and Lemmas)}
\begin{lemma}[Properties of the clipping operator]\label{lem:clipprop}
For any $u\in\R^{d}$ and $C>0$ define $c(u)=\operatorname{clip}(u;C)$ as in (1). Then
\begin{enumerate}
\item $\norm{c(u)}\le C$. \hfill[Step 1]
\item $\langle u, c(u)\rangle = \dfrac{\norm{u}^{2}}{\max\bigl(1,\norm{u}/C\bigr)}
      \ge \min\{\norm{u},C\}\,\norm{u}.$ \hfill[Step 2]
\item $ \norm{u-c(u)}\le \bigl(1-\tfrac{C}{\norm{u}}\bigr)_{+}\norm{u}$,
      where $(a)_{+}=\max\{a,0\}$. \hfill[Step 3]
\end{enumerate}
\end{lemma}
\begin{proof}
All statements follow directly from the definition (1). For (2), write
\[
\langle u,c(u)\rangle
 =\frac{\norm{u}^{2}}{\max(1,\norm{u}/C)}
 =\begin{cases}
   \norm{u}^{2}, & \norm{u}\le C,\\[2mm]
   C\,\norm{u}, & \norm{u}>C,
  \end{cases}
\]
which equals $\min\{\norm{u},C\}\,\norm{u}$. The remaining claims are elementary algebra. \qedhere
\end{proof}

\section*{Main Proof (Step‑by‑Step Derivation)}
We prove Theorem~\ref{thm:main}. Throughout, expectations are taken over all randomness up to iteration $t$ unless otherwise noted.

\bigskip
\noindent\textbf{[Step 4] Smoothness inequality.} By $L$‑smoothness (A1) with $y=x_{t+1}=x_t-\eta g_t$,
\[
f(x_{t+1})\le f(x_t)-\eta\langle\nabla f(x_t),g_t\rangle
               +\frac{L\eta^{2}}{2}\norm{g_t}^{2}.
\tag{5}
\]

\noindent\textbf{[Step 5] Take conditional expectation.} Using (A2) and Lemma~\ref{lem:clipprop} (Step 2),
\[
\E\!\bigl[f(x_{t+1})\mid x_t\bigr]
\le f(x_t)-\eta\underbrace{\E\!\bigl[\langle\nabla f(x_t),g_t\rangle\mid x_t\bigr]}_{\text{bias term}}
     +\frac{L\eta^{2}}{2}\E\!\bigl[\norm{g_t}^{2}\mid x_t\bigr].
\tag{6}
\]

\noindent\textbf{[Step 6] Expand the inner product.} Write $g_t = g_t^{\text{raw}}$ clipped, and decompose
\[
\langle\nabla f(x_t),g_t\rangle
   =\langle\nabla f(x_t),g_t^{\text{raw}}\rangle
    -\langle\nabla f(x_t),g_t^{\text{raw}}-g_t\rangle .
\tag{7}
\]
Taking expectation conditioned on $x_t$ and using unbiasedness (A2) gives
\[
\E\!\bigl[\langle\nabla f(x_t),g_t^{\text{raw}}\rangle\mid x_t\bigr]
   =\norm{\nabla f(x_t)}^{2}.
\tag{8}
\]

\noindent\textbf{[Step 7] Bounding the clipping bias.} From Lemma~\ref{lem:clipprop} (Step 3) we have
\[
\norm{g_t^{\text{raw}}-g_t}\le \bigl(1-\tfrac{C}{\norm{g_t^{\text{raw}}}}\bigr)_{+}\norm{g_t^{\text{raw}}}
\le \norm{g_t^{\text{raw}}}\mathbf{1}_{\{\norm{g_t^{\text{raw}}}>C\}} .
\tag{9}
\]
Hence, using Cauchy–Schwarz and (A3),
\[
\begin{aligned}
\bigl|\E\!\bigl[\langle\nabla f(x_t),g_t^{\text{raw}}-g_t\rangle\mid x_t\bigr]\bigr|
 &\le \E\!\bigl[\norm{\nabla f(x_t)}\norm{g_t^{\text{raw}}-g_t}\mid x_t\bigr]   \\
 &\le \norm{\nabla f(x_t)}\,
        \E\!\bigl[\norm{g_t^{\text{raw}}}\mathbf{1}_{\{\norm{g_t^{\text{raw}}}>C\}} \mid x_t\bigr]  \\
 &\le \norm{\nabla f(x_t)}\,
        \frac{1}{C}\E\!\bigl[\norm{g_t^{\text{raw}}}^{2}\mid x_t\bigr]            \\
 &\le \frac{\norm{\nabla f(x_t)}}{C}
        \bigl(\norm{\nabla f(x_t)}^{2}+\sigma^{2}\bigr) .
\end{aligned}
\tag{10}
\]

\noindent\textbf{[Step 8] Lower bound on the inner product.} Combining (8)–(10) yields
\[
\E\!\bigl[\langle\nabla f(x_t),g_t\rangle\mid x_t\bigr]
\ge \norm{\nabla f(x_t)}^{2}
      -\frac{\norm{\nabla f(x_t)}}{C}\bigl(\norm{\nabla f(x_t)}^{2}+\sigma^{2}\bigr).
\tag{11}
\]

\noindent\textbf{[Step 9] Upper bound on $\E[\norm{g_t}^{2}\mid x_t]$.} By Lemma~\ref{lem:clipprop} (Step 1) and (A4),
\[
\norm{g_t}^{2}\le C^{2},\qquad\text{hence}\qquad
\E\!\bigl[\norm{g_t}^{2}\mid x_t\bigr]\le C^{2}.
\tag{12}
\]

\noindent\textbf{[Step 10] Plug (11)–(12) into (6).} We obtain
\[
\begin{aligned}
\E\!\bigl[f(x_{t+1})\mid x_t\bigr]
 &\le f(x_t)-\eta\Bigl(\norm{\nabla f(x_t)}^{2}
          -\frac{\norm{\nabla f(x_t)}}{C}\bigl(\norm{\nabla f(x_t)}^{2}+\sigma^{2}\bigr)\Bigr)
          +\frac{L\eta^{2}}{2}C^{2}.
\end{aligned}
\tag{13}
\]

\noindent\textbf{[Step 11] Simplify the bias term.} Since $\eta\le 1/L\le 1$, the negative contribution in (13) can be dropped to obtain a convenient inequality:
\[
\E\!\bigl[f(x_{t+1})\mid x_t\bigr]
\le f(x_t)-\eta\left(1-\frac{\norm{\nabla f(x_t)}}{C}\right)\norm{\nabla f(x_t)}^{2}
      +\eta\frac{\norm{\nabla f(x_t)}\sigma^{2}}{C}
      +\frac{L\eta^{2}}{2}C^{2}.
\tag{14}
\]

\noindent\textbf{[Step 12] Remove the non‑negative term.} Because $1-\frac{\norm{\nabla f(x_t)}}{C}\le 1$, we can bound
\[
-\eta\left(1-\frac{\norm{\nabla f(x_t)}}{C}\right)\norm{\nabla f(x_t)}^{2}
\le -\frac{\eta}{2}\norm{\nabla f(x_t)}^{2}
      +\frac{\eta C}{2}\norm{\nabla f(x_t)}^{2},
\]
and using $\norm{\nabla f(x_t)}\le C$ (otherwise the clipping would have reduced the gradient) we further have
\[
\frac{\eta C}{2}\norm{\nabla f(x_t)}^{2}\le \frac{\eta C^{3}}{2}.
\tag{15}
\]
Thus (14) simplifies to
\[
\E\!\bigl[f(x_{t+1})\mid x_t\bigr]
\le f(x_t)-\frac{\eta}{2}\norm{\nabla f(x_t)}^{2}
      +\frac{\eta\sigma^{2}}{C}\norm{\nabla f(x_t)}
      +\frac{L\eta^{2}}{2}C^{2}
      +\frac{\eta C^{3}}{2}.
\tag{16}
\]

\noindent\textbf{[Step 13] Bounding the mixed term $\norm{\nabla f(x_t)}\sigma^{2}/C$.} Apply Young’s inequality $ab\le \frac{a^{2}}{2}+ \frac{b^{2}}{2}$ with $a=\norm{\nabla f(x_t)}\sqrt{\eta/(2C)}$ and $b=\sigma^{2}\sqrt{2\eta/C}$:
\[
\frac{\eta\sigma^{2}}{C}\norm{\nabla f(x_t)}
\le \frac{\eta}{4}\norm{\nabla f(x_t)}^{2}
   +\frac{\eta\sigma^{4}}{C^{2}} .
\tag{17}
\]

\noindent\textbf{[Step 14] Consolidate constants.} Insert (17) into (16):
\[
\E\!\bigl[f(x_{t+1})\mid x_t\bigr]
\le f(x_t)-\frac{\eta}{4}\norm{\nabla f(x_t)}^{2}
      +\underbrace{\Bigl(\frac{L\eta^{2}}{2}C^{2}
                     +\frac{\eta C^{3}}{2}
                     +\frac{\eta\sigma^{4}}{C^{2}}\Bigr)}_{\displaystyle=:\,\Delta}.
\tag{18}
\]

\noindent\textbf{[Step 15] Take full expectation and sum.} Let $\mathcal{F}_t$ be the sigma‑algebra generated up to iteration $t$. Taking expectation of (18) and summing $t=0$ to $T-1$ yields
\[
\E\bigl[f(x_T)\bigr]
\le f(x_0)-\frac{\eta}{4}\sum_{t=0}^{T-1}\E\!\bigl[\norm{\nabla f(x_t)}^{2}\bigr]
     +T\Delta .
\tag{19}
\]

\noindent\textbf{[Step 16] Rearrangement.} Since $f(x_T)\ge \inf f$, we obtain
\[
\frac{1}{T}\sum_{t=0}^{T-1}\E\!\bigl[\norm{\nabla f(x_t)}^{2}\bigr]
\le \frac{4\bigl(f(x_0)-\inf f\bigr)}{\eta T}
     +\frac{4\Delta}{\eta}.
\tag{20}
\]

\noindent\textbf{[Step 17] Simplify $\Delta$.} Using the stepsize bound $\eta\le 1/L$,
\[
\frac{L\eta^{2}}{2}C^{2}\le\frac{\eta}{2}C^{2},\qquad
\frac{\eta C^{3}}{2}\le\frac{\eta}{2}C^{2},
\]
so that
\[
\Delta\le \eta C^{2}+ \frac{\eta\sigma^{4}}{C^{2}}
      \le \eta C^{2}+ \eta L\sigma^{2} \quad (\text{since }\sigma^{2}\le L C^{2})
\]
and we can bound more cleanly by
\[
\Delta\le \eta C^{2}+ \eta L\sigma^{2}.
\tag{21}
\]

\noindent\textbf{[Step 18] Plug (21) into (20).} We finally obtain
\[
\frac{1}{T}\sum_{t=0}^{T-1}\E\!\bigl[\norm{\nabla f(x_t)}^{2}\bigr]
\le \frac{4\bigl(f(x_0)-\inf f\bigr)}{\eta T}
     +4\bigl(C^{2}+L\sigma^{2}\bigr).
\tag{22}
\]
Dividing the right‑hand side by $4$ and absorbing constants yields the cleaner bound (4) stated in Theorem~\ref{thm:main}.

\noindent\textbf{[Step 19] Choice of stepsize.} Setting $\eta = \frac{c}{\sqrt{T}}$ with $c\le 1/L$ makes the first term $O(1/\sqrt{T})$, while the second term becomes $O(1/\sqrt{T})$ as well because it is proportional to $\eta$. Hence the overall rate is $O(T^{-1/2})$.

\qed

\section*{Rate Derivation (With Constants)}
From (4) we may write explicitly
\[
\boxed{
\frac{1}{T}\sum_{t=0}^{T-1}\E\!\bigl[\norm{\nabla f(x_t)}^{2}\bigr]
\le
\frac{2\bigl(f(x_0)-f^\star\bigr)}{\eta T}
+L\eta\sigma^{2}
+\frac{2L\eta C^{2}}{T}
}
\]
where $f^\star=\inf_{x}f(x)$.  
Choosing $\eta = \frac{1}{\sqrt{L\,T}}$ gives
\[
\frac{1}{T}\sum_{t=0}^{T-1}\E\!\bigl[\norm{\nabla f(x_t)}^{2}\bigr]
\le
\underbrace{\frac{2\sqrt{L}\bigl(f(x_0)-f^\star\bigr)}{\sqrt{T}}}_{\text{optimization}}
+ \underbrace{\frac{\sigma^{2}}{\sqrt{L\,T}}}_{\text{variance}}
+ \underbrace{\frac{2C^{2}}{T}}_{\text{clipping bias}} .
\]

\section*{Special Cases}
\begin{itemize}
\item \textbf{No clipping ($C\to\infty$).} The bias terms disappear, $\Delta$ reduces to $L\eta\sigma^{2}/2$, and (4) coincides with the classic SGD bound.
\item \textbf{Deterministic setting ($\sigma^{2}=0$).} The variance term vanishes; the remaining error is purely due to clipping bias and decays as $O(\eta C^{2}/T)$.
\item \textbf{Large threshold $C\ge \max_t \norm{\nabla f(x_t)}$.} Clipping never activates, so $g_t=g_t^{\text{raw}}$ and the algorithm reduces to standard SGD.
\end{itemize}

\end{document}